\section{Identifying the homomorphism}\label{sec:identification}

We now turn the prime correspondence into equality of group homomorphisms. The argument uses one fixed coefficient prime $r$, although the auxiliary prime witnessing activity is allowed to vary. This distinction is what makes the correspondences constructed in Section~\ref{sec:matching} sufficient.

Retain the map $\beta:T=G_k\to P$, the infinite set $\mathcal L$, and the fields $B=\Q(i)$ and $K$ from Sections~\ref{sec:generation}--\ref{sec:matching}. Thus $G_K=\beta^{-1}(P_B)$, and Theorem~\ref{thm:normalized-map} gives
\begin{equation}\label{eq:fixed-r-log}
 \log\chi_r\circ\beta=\log\chi_{k,r}\qquad(r\in\mathcal L).
\end{equation}
Write $\sigma:T\to P$ for ordinary restriction.

\begin{proposition}\label{prop:identify-normalized}
There is an element $a\in P$ such that $\beta=\Ad(a)\circ\sigma$ on $T$.
\end{proposition}
\begin{proof}
Fix a finite Galois extension $L/\Q$ inside $\Q^s$ containing $B$, and put $R=\Gal(L/\Q)$. Choose a finite Galois extension $F'/\Q$ containing $K$ and $L$ such that $\beta(G_{F'})\subset P_L$. This is possible by taking the core in $\mathcal G$ of the appropriate open subgroup of $T$. Fix $r\in\mathcal L$.

The maps $\beta$ and $\sigma$, restricted to $G_{F'}$, give two pullbacks
\begin{equation}\label{eq:two-pullbacks}
 i_\beta,i_\sigma:\X rL\longrightarrow\X r{F'}.
\end{equation}
Both are injective: their underlying homomorphisms have open image, and an additive $\Q_r$-valued character vanishing on a finite-index subgroup is zero. For $a\in P$ let $\sigma_a=\Ad(a)\circ\sigma$ and let $i_a$ be its pullback. Since $L/\Q$ is Galois, $\sigma_a(G_{F'})\subset P_L$. The map $i_a$ depends only on $a$ modulo $P_L$, because inner automorphisms of $P_L$ act trivially on additive characters.

\emph{Frobenius evaluations.}
Let $q\ne r$ range over rational primes split completely in $F'$. By Proposition~\ref{prop:same-characteristic}, for all but finitely many such $q$ there is a prime $v$ of $K$ above $q$ whose decomposition group maps into a decomposition group at a prime $w$ of $B$ above the same $q$. Choose a prime $v'$ of $F'$ above $v$ and compatible prolongations. Then
\[
 \beta(D_{v'})\subset D_u\subset P_L
\]
for a prime $u$ of $L$ above $w$. Every prime here has absolute residue degree one, because $q$ splits completely in $F'$.

Choose a lift $h$ of arithmetic Frobenius in $D_{v'}$. Let $a_q\in\widehat{\Z}$ be the residue exponent of $\beta(h)$ at $u$. Equation~\eqref{eq:fixed-r-log} gives
\[
 (a_q)_r\log_r q=\log_r q.
\]
The number $\log_r q$ is nonzero: $q\in\Z_r^\times$ is a positive rational integer larger than one and is not a root of unity. Consequently $(a_q)_r=1$. Every $x\in\X rL$ is unramified at $u$, so
\begin{equation}\label{eq:matched-evaluation}
 i_\beta(x)(\Fr_{v'})=x(\Fr_u).
\end{equation}
This calculation uses only decomposition containment from the active correspondence. It does not require that its auxiliary prime be $r$.

Under ordinary restriction, the same Frobenius maps to Frobenius at another prime of $L$ above $q$. An element of $R$ carries this prime to $u$. Choose a lift $a\in P$ of that element. Then
\[
 i_\beta(x)(\Fr_{v'})=i_a(x)(\Fr_{v'})
 \qquad\text{for every }x\in\X rL.
\]
The choice of $a$ modulo $P_L$ may initially depend on $q$. Since $R$ is finite, an infinite set of distinct $q$ uses one fixed choice. For each $x$, the difference $i_\beta(x)-i_a(x)$ therefore vanishes at infinitely many degree-one primes of $F'$. Proposition~\ref{prop:character-detection} gives
\begin{equation}\label{eq:equal-pullbacks}
 i_\beta=i_a.
\end{equation}

\emph{From characters to the finite group.}
The field $F'/\Q$ is Galois, so $T$ acts on $\X r{F'}$ by conjugation. The two pullbacks intertwine this action with the respective actions on $\X rL$ induced by $\beta$ and $\sigma_a$. If $i$ denotes their common injective value in \eqref{eq:equal-pullbacks}, then for $t\in T$ and $x\in\X rL$,
\[
 i\bigl(\beta(t)\cdot x\bigr)
   =t\cdot i(x)
   =i\bigl(\sigma_a(t)\cdot x\bigr).
\]
The field $L$ is totally imaginary because it contains $B$. By Proposition~\ref{prop:signed-representation}, $R$ acts faithfully on $\X rL$. Hence
\[
 \beta(t)\bmod P_L=\sigma_a(t)\bmod P_L\qquad(t\in T).
\]

\emph{One conjugation for all finite quotients.}
For each such $L$, let
\[
 C_L=\{a\in P:\ \beta(t)\equiv\Ad(a)(\sigma(t))\pmod{P_L}
                    \text{ for every }t\in T\}.
\]
We have proved that $C_L$ is nonempty. It is closed, and indeed clopen, since the defining condition depends only on $a$ modulo $P_L$. If $L$ contains finitely many fields $L_1,\ldots,L_s$, then $C_L\subset\bigcap_j C_{L_j}$. Composita thus give the finite-intersection property. Compactness of $P$ supplies an element in every $C_L$. The extensions $L$ containing $B$ are cofinal among finite Galois subextensions of $\Q^s/\Q$, so this element gives the asserted equality in $P$.
\end{proof}

We remove the normalizer extension and return to the original source.
\begin{theorem}\label{thm:prosolvable-rigidity}
Let $F$ be a number field in $\overline{\Q}$. Every continuous open homomorphism $\varphi:G_F\to P$ is conjugate in $P$ to ordinary restriction.
\end{theorem}
\begin{proof}
As in Lemma~\ref{lem:normalizer-extension}, choose an open normal subgroup $U$ of $\mathcal G$ contained in $G_F$, and extend $\varphi|_U$ to the normalizer of its kernel. Proposition~\ref{prop:identify-normalized} shows that, after one conjugation of $\varphi$, it agrees on $U$ with ordinary restriction $\sigma:G_F\to P$.

For $g\in G_F$ and $u\in U$, normality of $U$ gives
\[
 \varphi(g)\sigma(u)\varphi(g)^{-1}
 =\varphi(gug^{-1})
 =\sigma(gug^{-1})
 =\sigma(g)\sigma(u)\sigma(g)^{-1}.
\]
Thus $\sigma(g)^{-1}\varphi(g)$ centralizes $\sigma(U)=\varphi(U)$, an open subgroup of $P$. Lemma~\ref{lem:centralizer} makes this element trivial. The equality holds for every $g\in G_F$.
\end{proof}

\begin{proof}[Proof of Theorem~\ref{thm:main}]
Choose copies of the given fields inside $\overline{\Q}$. By Lemma~\ref{lem:solvable-closure}, each $E_i$ contains $\Q^s$. Form the composite
\begin{equation}\label{eq:final-composite}
 G_{F_1}\longrightarrow G_1\xrightarrow{\alpha}G_2
     \longrightarrow P,
\end{equation}
where the outer arrows are restrictions. The first is surjective. The last has open image, since its composite with $G_{F_2}\twoheadrightarrow G_2$ is ordinary restriction from an open subgroup of $\mathcal G$ to $P$. Continuous surjective homomorphisms between profinite groups are open, so \eqref{eq:final-composite} is open.

By Theorem~\ref{thm:prosolvable-rigidity}, this composite is conjugate to ordinary restriction. Cyclotomic characters are unchanged by conjugation. All roots of unity lie in $\Q^s$, so the full cyclotomic characters factor through the groups in \eqref{eq:final-composite}. Surjectivity of the first arrow therefore yields
\[
 \chi_{E_2/F_2}\circ\alpha=\chi_{E_1/F_1}
 \quad\text{in }\widehat{\Z}^{\times}.
\]
Theorem~\ref{thm:hoshi} gives an equivariant embedding $E_2\hookrightarrow E_1$. Its uniqueness is Uchida's theorem cited after Theorem~\ref{thm:hoshi}.
\end{proof}
